% Mechanical relocation generated from the preservation ledger.
% Existing prose is included byte-for-byte from preserved blocks.

\section{Intrinsic Underdetermination of OD Estimation}
\label{sec:intrinsic-underdetermination}

OD estimation from passenger counts is an inverse problem.  Let
$\bm f\in\mathbb{R}_{+}^{n}$ collect the unknown OD-time demands and let
$\bm y\in\mathbb{N}^{m}$ collect the observed boarding and alighting counts.
For fixed network and route-choice assumptions, the expected observations can
be written schematically as
\[
  \mathbb{E}[\bm y\mid\bm f]=\bm\mu(\bm f).
\]
In the fixed-routing case this relationship is affine,
\begin{equation}
  \bm\mu(\bm f)=\rho\left(B\bm f+\bm b^{\mathrm{fix}}\right),
  \label{eq:part1-linear-observation}
\end{equation}
where a column of $B$ is the expected pattern of observed events generated by
one unit of demand in an OD-time cell.  This compact expression exposes the
central difficulty: passenger counts observe events along journeys, not the
journeys' origins and final destinations directly.

The problem is underdetermined whenever distinct feasible demands produce the
same expected counts.  In the linear case, if a nonzero direction $\bm v$
satisfies $B\bm v=\bm 0$, then $\bm f$ and $\bm f+\bm v$ are observationally
equivalent whenever both are feasible.  More generally, nearly collinear
columns of $B$ create weakly identified directions: the data can distinguish
them only through small differences that may be overwhelmed by measurement
noise.  Having more count records than unknowns is therefore neither necessary
nor sufficient for identification; the rank and conditioning of the response,
the observation coverage, and the imposed model restrictions matter.

Several common features aggravate this ambiguity.  OD pairs may share most of
their vehicle legs, transfer boardings are indistinguishable from initial
boardings in raw stop totals, and aggregation over stops or time can erase the
few events that distinguish two journeys.  Missing observations create no
constraint at all.  Even counts at every observed boarding and alighting event
need not uniquely identify every OD-time cell, because those events are linked
through unobserved passenger trajectories.

Consequently, an estimator cannot manufacture information absent from the
observations.  Maximum likelihood selects among demands using only the
likelihood and may be unstable or nonunique in weak directions.  MAP estimation
adds explicit prior or regularization assumptions and selects a stable point,
but stability must not be confused with empirical identification.  A reduced
demand model, such as gravity, constrains many OD cells through a small number
of shared parameters; this can make the parameter vector identifiable while
remaining conditional on stronger behavioral and production assumptions.
Bayesian inference expresses the remaining uncertainty, but its posterior is
also shaped by these assumptions where the likelihood is weak.

The report therefore distinguishes three questions throughout:
\begin{enumerate}
  \item \emph{structural feasibility}: can the timetable connect the OD-time
        pair under the accepted transfer and waiting rules?
  \item \emph{statistical identification}: do the available observations
        distinguish the retained parameters?
  \item \emph{model adequacy}: does the fitted model reproduce relevant count
        patterns without systematic residual structure?
\end{enumerate}
A structural zero answers only the first question.  Regularization or dimension
reduction addresses the second through assumptions.  Neither guarantees the
third, which must be assessed separately.
